\documentclass[11pt,reqno]{amsart}

\usepackage[utf8]{inputenc}
\usepackage[T1]{fontenc}
\usepackage{amsmath,amssymb,amsthm,mathtools}
\usepackage{mathrsfs}
\usepackage{geometry}
\geometry{a4paper, margin=1in}
\usepackage{hyperref}
\hypersetup{
    colorlinks=true,
    linkcolor=blue,
    citecolor=red,
    urlcolor=teal
}
\usepackage{tikz}
\usetikzlibrary{arrows.meta, calc}
\usepackage{microtype}
\usepackage{booktabs}
\usepackage{array}
\usepackage{tabularx}

% --- Teoremas e Estruturas ---
\newtheorem{theorem}{Theorem}[section]
\newtheorem{lemma}[theorem]{Lemma}
\newtheorem{proposition}[theorem]{Proposition}
\newtheorem{corollary}[theorem]{Corollary}
\newtheorem{conjecture}[theorem]{Conjecture}

\theoremstyle{definition}
\newtheorem{definition}[theorem]{Definition}
\newtheorem{remark}[theorem]{Remark}
\newtheorem{example}[theorem]{Example}
\newtheorem{assumption}[theorem]{Assumption}
\newtheorem{algorithm}[theorem]{Algorithm}

% --- Atalhos e Operadores ---
\providecommand{\R}{\mathbb{R}}
\newcommand{\N}{\mathbb{N}}
\newcommand{\C}{\mathbb{C}}
\providecommand{\Hyp}{\mathbb{H}}
\providecommand{\Sph}{\mathbb{S}}
\providecommand{\II}{\mathrm{I\!I}}
\newcommand{\op}{\mathrm{op}}
\providecommand{\dif}{\,\mathrm{d}}
\newcommand{\dist}{\mathrm{dist}}
\newcommand{\reach}{\mathrm{reach}}
\newcommand{\medial}{\mathcal{M}}
\providecommand{\Hsn}{\mathcal{H}}
\newcommand{\OmegaBar}{\bar{\Omega}}

\numberwithin{equation}{section}

\DeclareMathOperator*{\argmin}{argmin}
\providecommand{\Hyp}{\mathbb{H}}
\providecommand{\Sph}{\mathbb{S}}
\providecommand{\II}{\mathrm{II}}
\providecommand{\R}{\mathbb{R}}
\providecommand{\Area}{\mathrm{Area}}
\providecommand{\Hsn}{\mathcal{H}}
\providecommand{\BV}{\mathrm{BV}}
\providecommand{\TV}{\mathrm{TV}}
\providecommand{\cutnorm}[1]{\|#1\|_{\square}}

\begin{document}

\title[Minimax Curvature in Non-Euclidean Geometries and General Relativity]{Minimax Extrinsic Curvature in Non-Euclidean Geometries and General Relativity: Variational Theory, Space Forms, and ADM Spacetime Slicings}

\author{Reinaldo Maia Silva-Filho}
\address{Programa de P\'os-Gradua\c{c}\~ao em Estat\'istica e Experimenta\c{c}\~ao Agropecu\'aria (PPGEE/DES) \\
Universidade Federal de Lavras (UFLA), Lavras, MG, Brazil}
\email{reinaldo.filho1@estudante.ufla.br}
\thanks{This research was financed in part by the Coordena\c{c}\~ao de Aperfei\c{c}oamento de Pessoal de N\'ivel Superior -- Brasil (CAPES) -- Finance Code 001.}

\date{\today}

\begin{abstract}
We study the $L^\infty$ minimax extrinsic curvature problem of Chapter~7, minimizing $\|\II_M\|_{L^\infty(M)}$ over submanifolds of an obstacle domain with prescribed boundary, in Riemannian and Lorentzian ambient manifolds $(N^n, g)$: space forms and $4D$ spacetimes of General Relativity.

Our primary contributions are:
(1) \textbf{Non-Euclidean shape operator:} we set up the minimax problem with the classical Gauss--Codazzi--Ricci equations and prove the sectional--extrinsic coupling $K_M = \bar K + \kappa^2$ for umbilical hypersurfaces of a space form of curvature $\bar K$;
(2) \textbf{Space forms ($\Sph^n(+c^2)$ and $\Hyp^n(-c^2)$):} we compute the extrinsic curvature of standard examples (horospheres, equidistant hypersurfaces, tubes, small spheres, Clifford tori) and prove that the optimal U-turn curvature in a tube of width $w$ about a geodesic is $2/w$ in $\R^2$, $c\coth(cw/2)$ in $\Hyp^2(-c^2)$ and $c\cot(cw/2)$ in $\Sph^2(+c^2)$, attained by geodesic semicircles, so negative ambient curvature raises the required bending and positive curvature lowers it;
(3) \textbf{3+1 ADM slicings:} on slices of bounded extrinsic curvature the ADM constraints bound the scalar curvature, and the Raychaudhuri equation shows that maintaining zero expansion of a sheared geodesic congruence violates the strong energy condition;
(4) \textbf{Horizons, wormholes and junctions:} we compute the curvature of the Reissner--Nordstr\"om horizon in Painlev\'e--Gullstrand slicing, show that a Morris--Thorne throat is a minimal surface of its static slice, and, when the energy density at the throat is nonnegative, bound its null-energy violation by the throat's intrinsic Gauss curvature;
(5) \textbf{Relativistic trajectories:} we identify proper acceleration with the second fundamental form of a worldline and formulate the minimax problem for timelike curves in black-hole exteriors.
\end{abstract}

\maketitle

\tableofcontents
\newpage

\section{Introduction and Physical Motivations}

\subsection{From Euclidean Spaces to Curved Manifolds and Spacetime}
In Euclidean space $\R^n$, the second fundamental form $\II(X,Y) = (\bar{\nabla}_X Y)^\perp$ measures the deviation of a submanifold from flat affine planes. However, in physical geometry and General Relativity, the ambient manifold $(N^n, g)$ is inherently curved:
\begin{itemize}
    \item In \textbf{Riemannian Space Forms} ($\Sph^n$ of constant sectional curvature $+c^2$, and $\Hyp^n$ of curvature $-c^2$), submanifolds experience geometric drift induced by background curvature.
    \item In \textbf{Lorentzian Spacetimes} $(\mathcal{M}^{3+1}, g_{\mu\nu})$, the extrinsic curvature of spacelike Cauchy slices $\Sigma \subset \mathcal{M}$ (denoted by $K_{ij}$) represents the canonical momentum conjugate to the spatial metric $\gamma_{ij}$ in the Arnowitt-Deser-Misner (ADM) formulation of Einstein's field equations.
\end{itemize}

For membranes, spacetime slices or relativistic probes in strong gravitational fields (near black-hole throats, horizons or stellar surfaces), integral functionals such as area ($L^1$) or Willmore energy ($L^2$) control the curvature only on average. They do not bound its pointwise maximum, which is the relevant quantity when a local bound (a maximal stress, a maximal proper acceleration) must be respected.

This motivates the general \textbf{Non-Euclidean $L^\infty$ Minimax Extrinsic Curvature Problem}:
\begin{equation}
\label{eq:noneuclidean_problem}
\kappa^*_r(N, g, \Omega, \Sigma, V) := \inf_{M \in \mathcal{A}_r(N, g, \Omega, \Sigma, V)} \|\II_M\|_{L^\infty(M, g)}
\end{equation}
where, as in Chapter~7, $\mathcal{A}_r(N, g, \Omega, \Sigma, V)$ denotes the class of embedded $C^r$-submanifolds $M^k \subset \bar{\Omega} \subset N^n$ with boundary $\partial M = \Sigma^{k-1}$ and $k$-volume at most $V$, and $\mathcal{A}_{1,1}$ denotes the corresponding class of $C^{1,1}$ submanifolds. The norm $\|\II_M\|$ is that of Definition~\ref{def:opnorm}.

\subsection{Causal Immersions and the Prohibition of Self-Intersections in Spacetime}
\label{subsec:causality_self_intersection}
A profound physical and mathematical consequence arises when the ambient space is a Lorentzian spacetime $(\mathcal{M}^{3+1}, g_{\mu\nu})$. In contrast to Euclidean spaces where curves may self-intersect to relieve curvature, causal submanifolds in physically viable spacetimes are constrained by chronology:
\begin{enumerate}
    \item \textbf{Timelike Particle Trajectories ($k=1$):} For a massive observer with worldline $\gamma(\tau)$ parameterized by proper time $\tau$, the four-velocity satisfies $g(\dot{\gamma}, \dot{\gamma}) = -1$. The extrinsic curvature is the norm of the proper four-acceleration:
    \begin{equation}
    a^\mu = \nabla_u u^\mu \implies \|\II_\gamma\|_{\op, g} = \sqrt{g(a^\mu, a_\mu)}
    \end{equation}
    If $\gamma(\tau_1) = \gamma(\tau_2)$ with $\tau_1 < \tau_2$, the restriction $\gamma|_{[\tau_1,\tau_2]}$ is a closed timelike curve (CTC). \textbf{Self-intersections of timelike curves are therefore excluded in every chronological spacetime}, in particular in every stably causal one (the absence of CTCs is what Hawking's chronology protection conjecture \cite{hawking1992} asserts to be enforced dynamically). This applies to timelike curves only; spacelike curves in a Cauchy surface may self-intersect. Consequently, for timelike trajectories in chronological spacetimes the infimum of the peak acceleration over immersed curves equals the infimum over embedded curves: the immersion relaxation and the embedding problem coincide.
    \item \textbf{Spacelike ADM Slices ($k=3$):} Cauchy hypersurfaces $\Sigma_t \subset \mathcal{M}$ are achronal, so no two of their points are joined by a timelike curve, and they are embedded.
\end{enumerate}

\section{Geometric Foundations in Non-Euclidean Manifolds}

\subsection{Covariant Geometry of Submanifolds in \texorpdfstring{$(N^n, g)$}{(N\textasciicircum n, g)}}
Let $(N^n, g)$ be a smooth pseudo-Riemannian manifold of dimension $n$ with Levi-Civita connection $\bar{\nabla}$ and curvature tensor $\bar{R}$. Let $M^k \subset N$ be an immersed $C^r$-submanifold ($r \ge 2$) with induced metric $\gamma = i^* g$.

\begin{definition}[Second Fundamental Form and Non-Euclidean Shape Operator]
For tangent fields $X, Y \in \Gamma(TM)$ and normal field $\nu \in \Gamma(NM)$, the ambient covariant derivative decomposes into tangential and normal components via the Gauss-Weingarten formulas:
\begin{align}
\bar{\nabla}_X Y &= \nabla_X Y + \II(X, Y) \\
\bar{\nabla}_X \nu &= -A_\nu(X) + \nabla^\perp_X \nu
\end{align}
where $\nabla$ is the induced connection on $M$, $\nabla^\perp$ is the normal connection, $\II: TM \times TM \to NM$ is the symmetric bilinear \emph{second fundamental form}, and $A_\nu: TM \to TM$ is the \emph{shape operator} defined by:
\begin{equation}
g(A_\nu(X), Y) = g(\II(X, Y), \nu)
\end{equation}
\end{definition}

\begin{theorem}[Equations of Gauss, Codazzi-Mainardi, and Ricci; classical, see {\cite[Ch.~4]{oneill1983}}, {\cite[Ch.~2--3]{gourgoulhon2012}}]
With the curvature convention $R(X,Y) = \nabla_X\nabla_Y - \nabla_Y\nabla_X - \nabla_{[X,Y]}$, the ambient curvature tensor $\bar{R}$ and intrinsic curvature $R$ of $M^k$ satisfy:
\begin{align}
\label{eq:gauss}
g(R(X,Y)Z, W) &= g(\bar{R}(X,Y)Z, W) + g(\II(X,W), \II(Y,Z)) - g(\II(X,Z), \II(Y,W)) \\
\label{eq:codazzi}
(\bar{R}(X,Y)Z)^\perp &= (\bar{\nabla}_X \II)(Y,Z) - (\bar{\nabla}_Y \II)(X,Z) \\
\label{eq:ricci}
g(R^\perp(X,Y)\nu_1, \nu_2) &= g(\bar{R}(X,Y)\nu_1, \nu_2) + g([A_{\nu_1}, A_{\nu_2}]X, Y)
\end{align}
(O'Neill \cite{oneill1983} defines the curvature tensor with the opposite sign; the equations above are written in the convention stated here.)
\end{theorem}

\begin{definition}[Non-Euclidean Operator Norm]
\label{def:opnorm}
We use the operator norm in the following three situations, which are the ones that occur in this chapter:
\begin{enumerate}
    \item[(a)] $g$ is Riemannian, so that $T_pM$ and $N_pM$ are both positive definite (every submanifold of a Riemannian manifold; in Lorentzian signature a spacelike $T_pM$ always has a normal space containing a timelike vector, so this case does not occur there);
    \item[(b)] $T_pM$ is spacelike and $N_pM$ is a timelike line (spacelike hypersurfaces of a Lorentzian manifold);
    \item[(c)] $T_pM$ is a timelike line (timelike curves), so that $N_pM = T_pM^\perp$ is positive definite.
\end{enumerate}
In these cases the \emph{operator norm} of $\II_p$ is
\begin{equation}
\|\II_p\|_{\op, g} := \sup_{\substack{v \in T_pM \\ |g(v,v)| = 1}} \sqrt{|g(\II_p(v,v), \II_p(v,v))|} = \max_{\substack{\nu \in N_pM \\ |g(\nu,\nu)| = 1}} \ \sup_{\substack{v \in T_pM \\ |g(v,v)| = 1}} |g(\II_p(v,v), \nu)|,
\end{equation}
and in cases (a) and (b) the right-hand side equals $\max_\nu \max_i |\kappa_i^\nu(p)|$, the largest principal curvature over unit normals. The unit normals form a compact set in all three cases, so the maximum exists. For a spacelike surface of codimension two in a four-dimensional spacetime (a horizon cross-section, say) the normal space is Lorentzian and its unit vectors form a non-compact set. For the sphere $t = 0$, $r = R$ in Minkowski space and $\nu = \cosh\eta\, \partial_t + \sinh\eta\, \partial_r$, one finds $|g(\II(v,v),\nu)| = \sinh\eta/R \to \infty$ as $\eta \to \infty$, while $\sqrt{|g(\II(v,v),\II(v,v))|} = 1/R$. For such surfaces we do not use a spacetime norm; we measure the curvature of $M$ inside a chosen spacelike hypersurface $\Sigma \supset M$, as in Proposition~\ref{prop:rn_horizon}.
\end{definition}

\section{Exact Solutions in Space Forms: Spheres \texorpdfstring{$\Sph^n$}{S\textasciicircum n} and Hyperbolic Spaces \texorpdfstring{$\Hyp^n$}{H\textasciicircum n}}

\subsection{Curvature Coupling Theorem in Space Forms}
Let $N^n(\bar K)$ be an $n$-dimensional Riemannian space form of constant sectional curvature $\bar K \in \R$ (so $\bar K = +c^2$ for $\Sph^n(+c^2)$ and $\bar K = -c^2$ for $\Hyp^n(-c^2)$).
\begin{theorem}[Sectional-Extrinsic Coupling]
Let $n \ge 3$ and let $M^k \subset N^n(\bar K)$ be an umbilical hypersurface ($k = n-1 \ge 2$) of constant normal curvature $\kappa$. Then its intrinsic sectional curvature $K_M$ is constant and satisfies:
\begin{equation}
K_M = \bar K + \kappa^2
\end{equation}
\end{theorem}
\begin{proof}
Let $\nu$ be a unit normal. Umbilicity means $\II(X,Y) = \kappa\, g(X,Y)\,\nu$. In a space form $\bar{R}(X,Y)Z = \bar K\,(g(Y,Z)X - g(X,Z)Y)$. For orthonormal $X, Y \in T_pM$, the Gauss equation \eqref{eq:gauss} with $Z = Y$, $W = X$ gives
\[
K_M(X,Y) = g(R(X,Y)Y,X) = \bar K + g(\II(X,X),\II(Y,Y)) - g(\II(X,Y),\II(X,Y)) = \bar K + \kappa^2 - 0 .
\]
The value does not depend on the plane $X \wedge Y$, so $K_M$ is constant on $M$ whenever $\kappa$ is. (For $k = n-1 \ge 2$, the Codazzi equation gives $X(\kappa)\,Y - Y(\kappa)\,X = 0$ for independent tangent $X, Y$, so $\kappa$ is constant on connected umbilical hypersurfaces and the hypothesis is automatic.)
\end{proof}

\subsection{Constant-Curvature Examples in Hyperbolic Space \texorpdfstring{$\Hyp^n(-c^2)$}{H\textasciicircum n(-c\textasciicircum 2)}}
The following hypersurfaces of $\Hyp^n(-c^2)$ have explicit extrinsic curvature. They are examples, not solutions of a minimax problem with prescribed $(\Omega, \Sigma)$.
\begin{enumerate}
    \item \textbf{Horospheres ($\|\II\|_{\op} = c$):} In the Poincar\'e ball model, spheres tangent to the conformal boundary $\partial_\infty \Hyp^n$. Horospheres are umbilic with all principal curvatures equal to $c$, and their intrinsic metric is Euclidean ($K_M = -c^2 + c^2 = 0$).
    \item \textbf{Geodesic Hyperplanes ($\|\II\|_{\op} = 0$):} Totally geodesic hypersurfaces.
    \item \textbf{Equidistant Hypersurfaces ($\|\II\|_{\op} = c \tanh(c d)$):} The hypersurfaces at constant distance $d$ from a totally geodesic hyperplane. In the metric $\dif r^2 + \cosh^2(cr)\, g_{\Hyp^{n-1}}$ they are umbilic, with all principal curvatures equal to $c\tanh(cd)$.
    \item \textbf{Tubes about a geodesic ($\|\II\|_{\op} = c \coth(c d)$ for $n \ge 3$):} The hypersurfaces at constant distance $d$ from a geodesic. In the metric $\dif r^2 + c^{-2}\sinh^2(cr)\, \dif\Omega_{n-2}^2 + \cosh^2(cr)\, \dif t^2$ their principal curvatures are $c\coth(cd)$, with multiplicity $n-2$, and $c\tanh(cd)$. For $n = 2$ the tube is an equidistant curve with curvature $c\tanh(cd)$.
\end{enumerate}

\subsection{U-Turns in Corridors}
A \emph{U-turn} in a planar region is a $C^{1,1}$ arc-length parametrized curve $\gamma\colon[0,L] \to \R^2$ whose unit tangent at the end is the negative of its unit tangent at the start.

\begin{lemma}[Euclidean U-Turn Bound]
\label{lem:euclidean_uturn}
Let $S_w = \R \times [0,w]$ and let $\gamma \subset S_w$ be a U-turn starting in the direction $(1,0)$, with curvature $|\kappa| \le k$ a.e. Then $k \ge 2/w$. Equality is attained by a semicircle of diameter $w$.
\end{lemma}
\begin{proof}
Write $\gamma' = (\cos\theta, \sin\theta)$ with $\theta$ a continuous lift, $\theta(0) = 0$. Since $\gamma$ is $C^{1,1}$, $\theta$ is Lipschitz with $|\theta'| = |\kappa| \le k$ a.e. Because $\gamma'(L) = -\gamma'(0)$, $\theta(L) \in \pi + 2\pi\mathbb{Z}$, so $|\theta|$ attains $\pi$. Replacing $\gamma$ by its reflection $(x,y) \mapsto (x, w - y)$ if necessary, we may assume $\theta$ reaches $+\pi$ before $-\pi$. Let $s_b = \min\{s : \theta(s) = \pi\}$ and $s_a = \max\{s < s_b : \theta(s) = 0\}$. On $(s_a, s_b)$ we have $\theta \in (0, \pi)$, hence $\sin\theta > 0$, and by the intermediate value theorem every $t \in (0,\pi)$ is attained at least once. By the area formula for the Lipschitz map $\theta$,
\[
\begin{aligned}
w \;\ge\; y(s_b) - y(s_a) &= \int_{s_a}^{s_b} \sin\theta \dif s \;\ge\; \frac{1}{k}\int_{s_a}^{s_b} \sin\theta(s)\, |\theta'(s)| \dif s \\
&= \frac{1}{k}\int_0^\pi \sin t \; \#\{s \in (s_a,s_b) : \theta(s) = t\} \dif t \;\ge\; \frac{1}{k}\int_0^\pi \sin t \dif t = \frac{2}{k}.
\end{aligned}
\]
For the semicircle of radius $w/2$, $k = 2/w$ and the chain is an equality.
\end{proof}

In a space form the natural candidate is a geodesic circle. In $\Hyp^2(-c^2)$ a geodesic circle of radius $\rho$ has constant geodesic curvature $c\coth(c\rho)$, and in $\Sph^2(+c^2)$ it has $c\cot(c\rho)$ for $\rho < \pi/(2c)$. To compare tangent directions we use Fermi coordinates $(x, y)$ about a unit-speed geodesic $\sigma$, in which $y$ is the signed distance to $\sigma$ and
\begin{equation}
\label{eq:fermi}
g = \dif y^2 + G(y)^2 \dif x^2, \qquad G(y) = \begin{cases} \cosh(cy) & \text{in } \Hyp^2(-c^2), \\ \cos(cy) & \text{in } \Sph^2(+c^2), \\ 1 & \text{in } \R^2 . \end{cases}
\end{equation}
In $\Sph^2(+c^2)$ we take $|y| < \pi/(2c)$, where these coordinates are regular (then $x$ is periodic, which plays no role below). The fields $e_x = G(y)^{-1}\partial_x$ and $e_y = \partial_y$ form an orthonormal frame on the tube $T_w = \{|y| \le w/2\}$; $e_x$ is the parallel transport of $\sigma'$ along the geodesics orthogonal to $\sigma$. A \emph{U-turn} in $T_w$ is a $C^{1,1}$ arc-length parametrized curve $\gamma\colon [0,L] \to T_w$ with $\gamma'(0) = e_x$ and $\gamma'(L) = -e_x$. For $G \equiv 1$ this is the definition used in Lemma~\ref{lem:euclidean_uturn}.

\begin{theorem}[U-turns in tubes of space forms]
\label{prop:space_form_uturn}
Let $w > 0$, with $w < \pi/c$ in the case of $\Sph^2(+c^2)$. If $\gamma$ is a U-turn in $T_w$ whose geodesic curvature satisfies $|\kappa| \le k$ a.e., then
\begin{equation}
\label{eq:curved_uturn}
k \;\ge\; c \coth\left(\frac{c w}{2}\right) \ \text{in } \Hyp^2(-c^2), \qquad k \;\ge\; c \cot\left(\frac{c w}{2}\right) \ \text{in } \Sph^2(+c^2).
\end{equation}
Both bounds are attained by the half of the geodesic circle of radius $w/2$ centred on $\sigma$ that runs from $y = -w/2$ to $y = w/2$.
\end{theorem}
\begin{proof}
Write $\gamma' = \cos\theta\, e_x + \sin\theta\, e_y$ with $\theta$ a continuous lift, $\theta(0) = 0$; it exists because $(e_x, e_y)$ is a continuous frame on $T_w$. Let $N = -\sin\theta\, e_x + \cos\theta\, e_y$ be the left unit normal, so that $\bar\nabla_{\gamma'}\gamma' = \kappa N$, and write $y(s) = y(\gamma(s))$, so that $y' = \sin\theta$. The field $X = \partial_x$ is a Killing field, hence $g(\bar\nabla_{\gamma'}X, \gamma') = 0$ and
\begin{equation}
\label{eq:clairaut}
\frac{\dif}{\dif s}\, g(\gamma', X) = g(\kappa N, X) = -\kappa\, G(y) \sin\theta, \qquad g(\gamma', X) = G(y)\cos\theta .
\end{equation}
For $\kappa = 0$ this is Clairaut's relation. The function $u(s) = G(y(s))\cos\theta(s)$ is Lipschitz, and \eqref{eq:clairaut} holds a.e. As in the proof of Lemma~\ref{lem:euclidean_uturn}, $\theta(L) \in \pi + 2\pi\mathbb{Z}$; the reflection $y \mapsto -y$ is an isometry of $T_w$ that fixes $e_x$ and changes $\theta$ to $-\theta$, so we may assume that $\theta$ reaches $+\pi$ before $-\pi$. Let $s_b = \min\{s : \theta(s) = \pi\}$ and $s_a = \max\{s < s_b : \theta(s) = 0\}$. On $(s_a, s_b)$ we have $\sin\theta > 0$, so $y$ increases strictly from $a \coloneqq y(s_a) \ge -w/2$ to $b \coloneqq y(s_b) \le w/2$, and $G(y) > 0$ there. Integrating \eqref{eq:clairaut} over $[s_a, s_b]$,
\[
G(a) + G(b) = u(s_a) - u(s_b) = \int_{s_a}^{s_b} \kappa\, G(y) \sin\theta \dif s \;\le\; k \int_{s_a}^{s_b} G(y)\, y' \dif s = k \int_a^b G(y) \dif y .
\]
In $\Hyp^2(-c^2)$,
\[
\cosh(ca) + \cosh(cb) = 2\cosh\tfrac{c(a+b)}{2}\cosh\tfrac{c(b-a)}{2}, \qquad \int_a^b \cosh(cy)\dif y = \tfrac{2}{c}\cosh\tfrac{c(a+b)}{2}\sinh\tfrac{c(b-a)}{2},
\]
so $k \ge c\coth\frac{c(b-a)}{2} \ge c\coth\frac{cw}{2}$, since $0 < b - a \le w$ and $\coth$ is decreasing. In $\Sph^2(+c^2)$ the same identities with $\cos$ and $\sin$ give $k \ge c\cot\frac{c(b-a)}{2} \ge c\cot\frac{cw}{2}$, because $\cos\frac{c(a+b)}{2} > 0$ and $\cot$ is decreasing on $(0, \pi/2)$. For the geodesic semicircle, $\kappa$ is constant and equal to the right-hand side of \eqref{eq:curved_uturn}, $a = -w/2$ and $b = w/2$, so every inequality in the chain is an equality.
\end{proof}

With $G \equiv 1$ the same argument gives $2 \le k(b - a) \le kw$, a second proof of Lemma~\ref{lem:euclidean_uturn}. Since $x\coth x > 1 > x\cot x$ for $x \in (0,\pi/2)$, negative ambient curvature raises the optimal U-turn curvature above the Euclidean value $2/w$, and positive curvature lowers it.

\begin{remark}
The proof avoids comparing angles by parallel transport along $\gamma$, which in curved space would carry a holonomy term from Gauss--Bonnet. The conserved quantity of the Killing field $\partial_x$ absorbs that term. For $\Sph^2$ the argument uses only $G > 0$ on $T_w$, so it covers every $w < \pi/c$, including $w \ge \pi/(2c)$. The same argument gives, for any warped tube $\dif y^2 + G(y)^2\dif x^2$ with $G > 0$, the bound $k \ge \min_{-w/2 \le a < b \le w/2} (G(a) + G(b))/\int_a^b G$. The identity \eqref{eq:clairaut} and the bounds \eqref{eq:curved_uturn} were checked numerically, together with a search for U-turns with $0.98\,k$ that finds none (numerical check in the accompanying code repository).
\end{remark}

\subsection{Constant-Curvature Examples in Spheres \texorpdfstring{$\Sph^n(+c^2)$}{S\textasciicircum n(+c\textasciicircum 2)}}
In $\Sph^n(+c^2)$, totally geodesic submanifolds are great spheres $S^k$ ($\|\II\|_{\op} = 0$). Small spheres of radius $r$ have extrinsic curvature:
\begin{equation}
\|\II\|_{\op} = c \cot(c r), \qquad \lim_{r \to 0} c\cot(cr) = \frac{1}{r}.
\end{equation}
In $\Sph^3 \subset \R^4$, the Clifford torus $S^1(1/\sqrt{2}) \times S^1(1/\sqrt{2})$ has intrinsic curvature $K = 0$ and principal curvatures $\pm 1$, so $\|\II\|_{\op} = 1$.

\section{Applications to General Relativity and Spacetime Geometry}

\subsection{3+1 ADM Formalism and Minimax Cauchy Foliations}
In the $3+1$ ADM formulation of General Relativity, spacetime $(\mathcal{M}^4, g_{\mu\nu})$ is foliated by spacelike Cauchy hypersurfaces $\Sigma_t$ with future-directed timelike unit normal $n^\mu$ ($g(n, n) = -1$).
The spatial metric is $\gamma_{ij} = g_{ij} + n_i n_j$, and the \textbf{extrinsic curvature tensor} is:
\begin{equation}
\label{eq:Kij}
K_{ij} := -\frac{1}{2} \mathcal{L}_n \gamma_{ij} = -g(\bar{\nabla}_{\partial_i} n, \partial_j) = g(\bar{\nabla}_{\partial_i} \partial_j, n),
\end{equation}
where the last equality uses $g(n, \partial_j) = 0$. Since the normal component of a vector $V$ is $-g(V,n)\,n$, the second fundamental form of Definition~\ref{def:opnorm}(b) is $\II(\partial_i, \partial_j) = -K_{ij}\, n$, and $\|\II\|_{\op} = \max_i |\lambda_i|$ for the eigenvalues $\lambda_i$ of $K^i{}_j$. With this convention an expanding slicing has $K < 0$: for the hyperboloids $t^2 - |\mathbf{x}|^2 = \tau^2$ of Minkowski space, $K_{ij} = -\gamma_{ij}/\tau$.

\begin{theorem}[ADM Constraint Equations; classical, see \cite{adm1962}, {\cite[Ch.~3]{gourgoulhon2012}}]
The Einstein field equations $G_{\mu\nu} = 8\pi G T_{\mu\nu}$ project on $\Sigma_t$ as:
\begin{align}
\label{eq:hamiltonian_constraint}
\mathcal{H} &:= R^{(\gamma)} + K^2 - K_{ij}K^{ij} - 16\pi G \rho = 0 \quad &\text{(Hamiltonian Constraint)} \\
\label{eq:momentum_constraint}
\mathcal{M}^i &:= D_j (K^{ij} - \gamma^{ij} K) - 8\pi G J^i = 0 \quad &\text{(Momentum Constraint)}
\end{align}
where $K = \gamma^{ij} K_{ij}$ is the mean curvature, $\rho = T_{\mu\nu}n^\mu n^\nu$, $J^i = -\gamma^{i\mu} T_{\mu\nu} n^\nu$, and $\sigma_{ij} = K_{ij} - \frac{1}{3} K \gamma_{ij}$ is the trace-free part of $K_{ij}$.
\end{theorem}
\begin{proof}[Sketch]
Contracting the Gauss equation \eqref{eq:gauss} twice over $T\Sigma_t$ gives $2 G_{\mu\nu}n^\mu n^\nu = R^{(\gamma)} + K^2 - K_{ij}K^{ij}$. Contracting the Codazzi equation \eqref{eq:codazzi} once gives $\gamma^{i\mu} G_{\mu\nu} n^\nu = D_j K^{ij} - D^i K$, up to the sign convention for $K_{ij}$. Inserting $G_{\mu\nu} = 8\pi G T_{\mu\nu}$ yields both constraints. See \cite[\S 3.4]{gourgoulhon2012} for details.
\end{proof}

\begin{corollary}[Pointwise Bounds on Slices of Bounded Extrinsic Curvature]
\label{cor:adm_bounds}
If $\|K\|_{\op} \le \kappa$ at a point of $\Sigma_t$, then $K_{ij}K^{ij} \le 3\kappa^2$, $\sigma_{ij}\sigma^{ij} \le 3\kappa^2 - K^2/3$, and
\[
16\pi G\rho - 6\kappa^2 \;\le\; R^{(\gamma)} \;\le\; 16\pi G \rho + 2\kappa^2 .
\]
All bounds are sharp.
\end{corollary}
\begin{proof}
With principal curvatures $\lambda_i \in [-\kappa,\kappa]$ we have $K_{ij}K^{ij} = \sum\lambda_i^2 \le 3\kappa^2$ and $\sigma_{ij}\sigma^{ij} = K_{ij}K^{ij} - K^2/3$. The Hamiltonian constraint gives $R^{(\gamma)} - 16\pi G\rho = K_{ij}K^{ij} - K^2 = -2e_2(\lambda)$, where $e_2 = \lambda_1\lambda_2 + \lambda_2\lambda_3 + \lambda_3\lambda_1$. Since $e_2$ is affine in each $\lambda_i$ separately, its extreme values on the cube $[-\kappa,\kappa]^3$ are attained at vertices, where it takes the values $3\kappa^2$ (all signs equal) and $-\kappa^2$ (one sign different). Hence $e_2$ ranges exactly over $[-\kappa^2, 3\kappa^2]$, which gives $-6\kappa^2 \le R^{(\gamma)} - 16\pi G\rho \le 2\kappa^2$, with both ends attained. The bound $K_{ij}K^{ij} \le 3\kappa^2$ is attained at $\lambda_i \equiv \kappa$, and the bound on $\sigma_{ij}\sigma^{ij}$ is attained whenever all $|\lambda_i| = \kappa$.
\end{proof}

\begin{proposition}[Raychaudhuri Obstruction to Frozen Expansion]
\label{prop:raychaudhuri}
Let $V^\mu$ be the tangent field of a timelike geodesic congruence with zero vorticity, $\omega_{\mu\nu} = 0$ (for instance, the unit normals of a foliation whose normal observers are geodesic, i.e.\ with lapse $N \equiv 1$). If the expansion vanishes identically, $\theta \equiv 0$, on an open set $U$, then on $U$
\[
R_{\mu\nu}V^\mu V^\nu = -\sigma_{\mu\nu}\sigma^{\mu\nu} \le 0,
\]
with strict inequality wherever the shear is non-zero. In that case the Strong Energy Condition $R_{\mu\nu}V^\mu V^\nu \ge 0$ fails at every point of $U$ with $\sigma \ne 0$.
\end{proposition}
\begin{proof}
The Raychaudhuri equation \cite{raychaudhuri1955} for a geodesic congruence reads
$\dot\theta = -\tfrac{1}{3}\theta^2 - \sigma_{\mu\nu}\sigma^{\mu\nu} + \omega_{\mu\nu}\omega^{\mu\nu} - R_{\mu\nu}V^\mu V^\nu$.
With $\theta \equiv 0$ and $\omega = 0$ this reduces to $R_{\mu\nu}V^\mu V^\nu = -\sigma_{\mu\nu}\sigma^{\mu\nu}$. The shear is spatial, so $\sigma_{\mu\nu}\sigma^{\mu\nu} \ge 0$ with equality iff $\sigma = 0$.
\end{proof}

\begin{remark}[Posing a minimax slicing problem]
\label{rem:minimax_slicing}
The problem $\inf_\Sigma \operatorname*{ess\,sup}_{p \in \Sigma} \|K_p\|_{\op}$ over Cauchy hypersurfaces is degenerate without boundary conditions. In the Kruskal extension of Schwarzschild the slice $t = \text{const}$ (the Einstein--Rosen bridge) is time-symmetric, $K_{ij} \equiv 0$, so the infimum is $0$ and is attained by a slice that never enters the black-hole region. A meaningful problem has to prescribe how the slices behave, for instance that they coincide with a given slice near spatial infinity and cross the future event horizon, or that they reach a given surface in the black-hole interior. Proposition~\ref{prop:raychaudhuri} does not restrict such a problem: it concerns geodesic normals ($N = N(t)$) with $\theta \equiv 0$, a configuration that a minimizer of $\|K\|_{L^\infty}$ need not have. A constraint $R_{\mu\nu}V^\mu V^\nu \ge 0$ on the unit normal $V$ of the slices restricts the slicing only where the strong energy condition fails; where it holds (for instance in vacuum) the constraint is empty. We do not formulate or solve a constrained version here.
\end{remark}

\subsection{Apparent Horizons and Trapped Surfaces}
In black hole spacetimes, let $\mathcal{S} \subset \Sigma$ be a closed 2-surface ($k=2, n=4$). Let $l^\mu$ and $k^\mu$ be outgoing and ingoing future-directed null normal vector fields, normalized by $g(l, k) = -2$.

\begin{definition}[Null Expansion and Apparent Horizons]
The null expansions are
\begin{equation}
\theta_l := q^{ab}\, g(\bar{\nabla}_{e_a} l, e_b) = -q^{ab}\, g(\II(e_a, e_b), l), \qquad \theta_k := q^{ab}\, g(\bar{\nabla}_{e_a} k, e_b) = -q^{ab}\, g(\II(e_a, e_b), k),
\end{equation}
where $q_{ab}$ is the induced metric on $\mathcal{S}$, $(e_a)$ is a tangent frame, and the second equalities follow from $g(l, e_b) = g(k, e_b) = 0$. For the sphere $t = 0$, $r = R$ in Minkowski space with $l = \partial_t + \partial_r$ this gives $\theta_l = 2/R > 0$, as it should for an untrapped surface. A surface $\mathcal{S}$ is:
\begin{itemize}
    \item \textbf{Trapped:} if $\theta_l < 0$ and $\theta_k < 0$;
    \item \textbf{Marginally Outer Trapped Surface (MOTS):} if $\theta_l = 0$;
    \item \textbf{Apparent Horizon:} the outermost boundary of the region containing trapped surfaces.
\end{itemize}
\end{definition}

\begin{proposition}[Horizon Curvature in Painlev\'e--Gullstrand Slicing]
\label{prop:rn_horizon}
Let $(\mathcal{M}, g)$ be the Reissner--Nordstr\"om spacetime with $|Q| \le M$ (in particular Schwarzschild, $Q = 0$), $f(r) = 1 - 2M/r + Q^2/r^2$ and $r_+ = M + \sqrt{M^2 - Q^2}$. Write it in Painlev\'e--Gullstrand form
\[
\dif s^2 = -\dif T^2 + \big(\dif r + \beta(r)\, \dif T\big)^2 + r^2 \dif\Omega^2, \qquad \beta = \sqrt{2M/r - Q^2/r^2},
\]
valid for $r > Q^2/(2M)$, a range that contains $r_+$. Then:
\begin{enumerate}
    \item the slices $T = \text{const}$ are intrinsically flat, and the horizon sphere $\mathcal{S}_+ = \{r = r_+\}$ has both principal curvatures in the slice equal to $1/r_+$, so $\|\II_{\mathcal{S}_+ \subset \Sigma_T}\|_{\op} = 1/r_+$;
    \item $\mathcal{S}_+$ is a MOTS, $\theta_l = 0$.
\end{enumerate}
\end{proposition}
\begin{proof}
Expanding the PG line element gives $g_{TT} = -(1 - \beta^2) = -f$, $g_{Tr} = \beta$, $g_{rr} = 1$, so it is the Reissner--Nordstr\"om metric in the coordinate $T = t + \int \beta f^{-1}\dif r$. The induced metric on $T = \text{const}$ is $\dif r^2 + r^2 \dif\Omega^2$, which is flat. A round sphere of radius $r_+$ in flat space has principal curvatures $1/r_+$. The lapse is $N = 1$ and the shift is $\beta^r = \beta$. Since $\partial_T \gamma_{ij} = 0$, the convention \eqref{eq:Kij}, $K_{ij} = -\frac{1}{2N}(\partial_T\gamma_{ij} - D_i\beta_j - D_j\beta_i)$, gives $K_{ij} = D_{(i}\beta_{j)}$, hence $K^r{}_r = \beta'$ and $K^\theta{}_\theta = K^\phi{}_\phi = \beta/r$. With outward unit normal $s = \partial_r$ in the slice and $l = n + s$, the outgoing null expansion is $\theta_l = D_i s^i - K + K_{ij}s^is^j = 2/r - (\beta' + 2\beta/r) + \beta' = 2(1 - \beta)/r$. It vanishes at $r = r_+$, where $\beta(r_+)^2 = 1 - f(r_+) = 1$. (Symbolic check in the accompanying code repository.)
\end{proof}

\begin{remark}[Slicing dependence and rotation]
The value $1/r_+$ depends on the slicing. The static slices $t = \text{const}$ of the maximally extended spacetime all meet at the bifurcation sphere, which is a minimal surface of each of them: the curvature $\sqrt{f(r)}/r$ of the spheres $r = \text{const}$ in a static slice tends to $0$ as $r \to r_+$. For Kerr--Newman with $a \ne 0$ the horizon is not round, and no single number $1/r_+$ describes its curvature. Whether apparent horizons minimize $\|\II\|_{\op}$ among surfaces in some natural class is an open question that we do not address.
\end{remark}

\subsection{Israel Junction Conditions and Thin Shells}
Consider two spacetime manifolds $(\mathcal{M}^+, g^+)$ and $(\mathcal{M}^-, g^-)$ joined at a timelike hypersurface $\Sigma$ ($k=3$).

\begin{theorem}[Israel Junction Theorem; classical, \cite{israel1966}, see also \cite{mars1993}]
Let $[g_{ab}] = 0$ (metric continuous across $\Sigma$), let $n$ be the unit (spacelike) normal to $\Sigma$ pointing from $\mathcal{M}^-$ to $\mathcal{M}^+$, and in this subsection let $K_{ab} := h_a{}^\mu h_b{}^\nu \nabla_\mu n_\nu$, the convention of Israel. The Einstein field equations across $\Sigma$ in the sense of distributions produce a surface stress-energy tensor $S_{ab}$ given by the jump in the extrinsic curvature:
\begin{equation}
\label{eq:israel}
S_{ab} = -\frac{1}{8\pi G} \left( [K_{ab}] - h_{ab} [K] \right)
\end{equation}
where $[K_{ab}] = K_{ab}^+ - K_{ab}^-$ and $h_{ab}$ is the induced metric on $\Sigma$. With the convention \eqref{eq:Kij} of Section~4.1, $K_{ab} = -h_a{}^\mu h_b{}^\nu \nabla_\mu n_\nu$, and the sign of the right-hand side of \eqref{eq:israel} is reversed.
\end{theorem}

\begin{remark}[Relation to $C^{1,1}$ Regularity]
Two different objects must be kept apart. In the Israel setting the hypersurface $\Sigma$ is smooth and the \emph{spacetime metric} is only Lipschitz ($C^{0,1}$) across it, with a jump in its first normal derivatives, i.e.\ in $K_{ab}$. In Chapter~7 \cite{silvafilho2026minimax} the ambient metric is smooth and the admissible \emph{submanifolds} are $C^{1,1}$, so their second fundamental form is bounded but may jump. Minimizers of Dubins type are $C^{1,1}$ and not $C^2$; that $C^{1,1}$ is the optimal regularity of minimizers in general is expected there, not proved. The two situations share the feature that curvature is $L^\infty$ with jump discontinuities, and the same distributional calculus applies to both. They are not the same statement, and neither regularity result implies the other.
\end{remark}

\subsection{Traversable Wormholes and Exotic Matter Minimization}
In a Morris-Thorne traversable wormhole with metric:
\begin{equation}
ds^2 = -e^{2\Phi(r)} dt^2 + \frac{dr^2}{1 - b(r)/r} + r^2 (d\theta^2 + \sin^2\theta \, d\phi^2)
\end{equation}
The throat is located at $r = r_0$, where $b(r_0) = r_0$ and $b'(r_0) \le 1$.

\begin{proposition}[Throat Geometry and the Null Energy Bound]
\label{prop:wormhole_throat}
Assume $\Phi$ is $C^1$ near the throat, $b(r_0) = r_0$ and the flare-out condition $b'(r_0) < 1$ holds. Let $\Sigma = \{t = \text{const}\}$ and $\mathcal{S}_0 = \{r = r_0\} \subset \Sigma$. Then:
\begin{enumerate}
    \item $\Sigma$ is totally geodesic in spacetime ($K_{ij} = 0$), since the metric is static.
    \item The spheres $\{r = \text{const}\} \subset \Sigma$ have principal curvatures $\sqrt{1 - b(r)/r}\,/\,r$. At the throat these vanish, so $\mathcal{S}_0$ is a minimal surface of $\Sigma$, while its intrinsic Gauss curvature is $K_{\mathrm{Gauss}}(\mathcal{S}_0) = 1/r_0^2$.
    \item For the radial null vector $k$ normalized by $k^{\hat t} = k^{\hat r} = 1$ in the static orthonormal frame,
    \begin{equation}
        T_{\mu\nu}k^\mu k^\nu\big|_{r_0} = \rho + p_r = -\frac{1 - b'(r_0)}{8\pi G\, r_0^2} < 0,
    \end{equation}
    so the NEC is violated at the throat. If in addition $\rho(r_0) \ge 0$, equivalently $b'(r_0) \ge 0$, then
    \begin{equation}
        T_{\mu\nu}k^\mu k^\nu\big|_{r_0} \;\ge\; -\frac{1}{8\pi G}\, K_{\mathrm{Gauss}}(\mathcal{S}_0),
    \end{equation}
    with equality iff $b'(r_0) = 0$.
\end{enumerate}
\end{proposition}
\begin{proof}
(1) $\partial_t$ is a hypersurface-orthogonal Killing field, so $\mathcal{L}_n \gamma_{ij} = 0$. (2) In the induced metric $\dif r^2/(1 - b/r) + r^2\dif\Omega^2$ the unit normal to $r = \text{const}$ is $\nu = \sqrt{1 - b/r}\,\partial_r$, and $\tfrac{1}{2}\mathcal{L}_\nu(r^2\dif\Omega^2) = r\,\nu(r)\,\dif\Omega^2$, so the principal curvatures are $\pm\nu(r)/r = \pm\sqrt{1 - b/r}/r$ (the sign depends on the orientation), which vanish at $r_0$. The induced metric on $\mathcal{S}_0$ is that of a round sphere of radius $r_0$. (3) The Einstein tensor of the Morris--Thorne metric gives $8\pi G\rho = b'/r^2$ and $8\pi G p_r = -b/r^3 + 2(1 - b/r)\Phi'/r$ \cite{morristhorne1988}. At $r_0$ the $\Phi'$ term vanishes, so $8\pi G(\rho + p_r) = (b'(r_0) - 1)/r_0^2$. If $0 \le b'(r_0) < 1$, then $1 - b'(r_0) \le 1$, which is the stated inequality. Without the hypothesis $b'(r_0) \ge 0$ no such bound holds, since $1 - b'(r_0)$ is unbounded as $b'(r_0) \to -\infty$. (Symbolic check in the accompanying code repository.)
\end{proof}

\begin{remark}
By item (2), the throat has vanishing extrinsic curvature in its static slice; $1/r_0$ is the square root of its \emph{intrinsic} curvature, not an extrinsic curvature $K^\theta{}_\theta$. Whether an $L^\infty$-minimax construction as in Chapter~7 minimizes the total NEC violation among wormhole geometries with a given throat is an open question.
\end{remark}

\subsection{Relativistic Trajectory Optimization with Bounded Proper Acceleration}
Let $\gamma(\tau)$ be a future-directed timelike curve ($k=1, n=4$) representing the trajectory of a spacecraft in a curved spacetime $(\mathcal{M}, g)$.

\begin{proposition}[Proper Acceleration as Second Fundamental Form]
\label{prop:proper_accel}
Let $\gamma$ be a $C^2$ timelike curve parametrized by proper time, with $u = \dot\gamma$, $g(u,u) = -1$. Then its second fundamental form is $\II(u,u) = a \coloneqq \bar\nabla_u u$, the acceleration $a$ is spacelike or zero, and $\|\II_\gamma\|_{\op,g} = |a|_g = \sqrt{g(a,a)}$.
\end{proposition}
\begin{proof}
Differentiating $g(u,u) = -1$ along $\gamma$ gives $g(a,u) = 0$, so $a$ lies in the normal space $u^\perp$, which is spacelike. The tangential projection of $\bar\nabla_u u$ is $g(\bar\nabla_u u, u)\, u / g(u,u) = 0$, hence $\bar\nabla_u u = \II(u,u)$. Since $T\gamma$ is one-dimensional and spanned by the unit vector $u$, the operator norm is $|\II(u,u)|_g = |a|_g$.
\end{proof}

\begin{definition}[Minimax Worldline Problem]
\label{def:minimax_worldline}
Let $\Omega \subset \mathcal{M}$ be the exterior region of a black hole, and let $p, q \in \Omega$ with $q \in I^+(p)$. The \emph{minimax worldline problem} is
\begin{equation}
\kappa^*_{\mathrm{wl}}(\Omega; p, q) \coloneqq \inf_{\gamma} \operatorname*{ess\,sup}_{\tau} |a(\tau)|_g,
\end{equation}
where the infimum is taken over $C^{1,1}$ future-directed timelike curves in $\Omega$ from $p$ to $q$. By Proposition~\ref{prop:proper_accel} and Definition~\ref{def:opnorm}(c) this is an analogue of \eqref{eq:noneuclidean_problem} for $k = 1$, with the class $\mathcal{A}_{1,1}$, boundary $\{p, q\}$ and no volume bound. Physically, it minimizes the peak proper acceleration felt on board.
\end{definition}

\subsection{Winding Classes and the Relativistic Slingshot}
\label{subsec:relativistic_slingshot}
The global search of Chapter~9 enumerates homotopy classes of paths around obstacles. In a spatial slice of a black-hole spacetime, the natural obstacles are the horizons, and the relevant topology depends on dimension.

\begin{proposition}[Topology of Multi-Black-Hole Slices]
\label{prop:bh_topology}
Let $B_1, \dots, B_m \subset \R^3$ be disjoint closed balls and $\Sigma = \R^3 \setminus \bigcup_i B_i$. Then $\Sigma$ is simply connected, $\pi_1(\Sigma) = 1$. By contrast, if $D_1, \dots, D_m \subset \R^2$ are disjoint closed discs, then $\pi_1(\R^2 \setminus \bigcup_i D_i) \cong \mathbb{F}_m$, the free group on $m$ generators.
\end{proposition}
\begin{proof}
$\R^3 \setminus \bigcup_i B_i$ is homotopy equivalent to a wedge of $m$ two-spheres, obtained by shrinking each ball to a point and applying van Kampen along separating planes. Its fundamental group is the free product of $m$ copies of $\pi_1(S^2) = 1$. In the plane, $\R^2 \setminus \bigcup_i D_i$ is homotopy equivalent to a wedge of $m$ circles, whose fundamental group is $\mathbb{F}_m$.
\end{proof}

Winding numbers about horizons are therefore homotopy invariants only for motion confined to a two-dimensional surface, such as the equatorial plane of a static black hole or a (2+1)-dimensional model. For unconstrained motion in three spatial dimensions, every loop around a horizon can be contracted by passing over the horizon's pole.

\begin{proposition}[Causal Constraints on Minimax Worldlines]
\label{prop:causal_constraints}
Let $(\mathcal{M}, g)$ be a strongly causal, asymptotically flat black-hole spacetime with black-hole region $\mathcal{B} = \mathcal{M} \setminus J^-(\mathscr{I}^+)$. Items (2) and (3) concern the worldline problem posed on all of $\mathcal{M}$ rather than on the exterior region of Definition~\ref{def:minimax_worldline}.
\begin{enumerate}
    \item Every future-directed timelike curve from $p$ to $q$ lies in $J^+(p) \cap J^-(q)$, its interior points lie in $I^+(p) \cap I^-(q)$, and it is injective.
    \item If $p \in \mathcal{B}$ and $q \notin \mathcal{B}$, there is no future-directed causal curve from $p$ to $q$. The admissible class in the minimax worldline problem is then empty, and by convention $\kappa^*_{\mathrm{wl}} = +\infty$.
    \item If $p \notin \mathcal{B}$ and $q \in \mathcal{B} \cap I^+(p)$, the admissible class is non-empty and $\kappa^*_{\mathrm{wl}} < \infty$. Entering the black-hole region imposes no divergence of $\kappa^*_{\mathrm{wl}}$.
\end{enumerate}
\end{proposition}
\begin{proof}
(1) This follows from the definitions of $I^\pm$ and $J^\pm$. Injectivity holds because strong causality excludes closed causal curves. (2) If a causal curve went from $p$ to $q \in J^-(\mathscr{I}^+)$, transitivity would give $p \in J^-(\mathscr{I}^+)$, contradicting $p \in \mathcal{B}$. (3) Since $q \in I^+(p)$, there is a piecewise smooth future-directed timelike curve from $p$ to $q$; rounding its finitely many corners keeps it timelike and makes it $C^2$, and a $C^2$ curve on a compact parameter interval has bounded acceleration.
\end{proof}

\begin{conjecture}[Planar Relativistic Slingshot]
\label{conj:slingshot}
In the equatorial plane of Schwarzschild, fix events $p, q$ outside the photon sphere $r = 3M$ with $q \in I^+(p)$, and restrict to worldlines confined to the plane and to the region $r > 2M$. Then the minimax problem has a minimizer in every winding class $W$ about the horizon that contains an admissible worldline. Moreover, there are data $(p,q)$ for which the minimum over $W = \pm 1$ is strictly smaller than the minimum over $W = 0$.
\end{conjecture}

\begin{remark}
Only finitely many winding classes contain admissible worldlines. Along a future-directed timelike curve in the equatorial plane, $-f\dot t^2 + f^{-1}\dot r^2 + r^2\dot\phi^2 < 0$ with $f = 1 - 2M/r$, so
\[
\left|\frac{\dif\phi}{\dif t}\right| < \frac{\sqrt{f(r)}}{r} \le \max_{r > 2M}\frac{\sqrt{f(r)}}{r} = \frac{1}{3\sqrt3\, M},
\]
the maximum being attained at $r = 3M$. If $\Delta t$ is the Schwarzschild-time separation of $p$ and $q$ and $\phi_q - \phi_p$ is taken in $(-\pi, \pi]$, a worldline of winding number $W$ has $|\Delta\phi| = |\phi_q - \phi_p + 2\pi W| < \Delta t/(3\sqrt3 M)$, hence $|W| < \Delta t/(6\sqrt3\,\pi M) + \tfrac12$. For $\Delta t = 100M$ this gives $|W| \le 3$. A minimizer therefore cannot exist in every class $W \in \mathbb{Z}$, since all other classes are empty.
\end{remark}

\begin{remark}
Two further tempting statements are not correct. First, $\sup|a|$ need not diverge as the pericenter approaches $3M$: the effective potential is finite there. Second, an estimate of the form $\kappa^*_{W=\pm1} \approx M/(r_0^2\sqrt{1 - 3M/r_0})$ does not hold: circular timelike geodesics exist for $r_0 > 3M$ and require no thrust at all, so the orbital portion of a winding manoeuvre can have $a = 0$, and the cost comes entirely from matching the boundary data. Quantifying that cost is part of Conjecture~\ref{conj:slingshot}.
\end{remark}

\begin{remark}[Regularity of the optimal worldline]
\label{rem:worldline_regularity}
In the planar Euclidean problem, minimizers are Dubins-type concatenations of arcs and segments, which are $C^{1,1}$ and not $C^2$ (Chapter~7). If optimal worldlines have the same structure, their acceleration jumps and the jerk $\dot a$ is a measure. Whether the infimum over $C^2$ worldlines equals the infimum over $C^{1,1}$ worldlines is the Lorentzian analogue, for timelike curves, of the regularity-invariance conjecture of Chapter~7; Conjecture~\ref{thm:noneuclidean_invariance} below states it only for Riemannian ambient manifolds, and we have no proof for worldlines. If the equality holds and the infimum is not attained in $C^2$, physically realizable trajectories, which require bounded jerk, approach $\kappa^*_{\mathrm{wl}}$ without reaching it.
\end{remark}

\subsection{Singularity-Avoiding ADM Cauchy Slices via Bona-Mass\'o Hyperbolic Gauges}
\label{subsec:singularity_avoiding_adm}
In the Arnowitt-Deser-Misner (ADM) $3+1$ formulation of General Relativity, spatial hypersurfaces $\Sigma_t \subset \mathcal{M}$ are evolved according to the extrinsic curvature $K_{ij} = -\frac{1}{2\alpha} (\partial_t - \mathcal{L}_\beta) \gamma_{ij}$. In Schwarzschild, maximal slicing ($K = \operatorname{tr} K = 0$) avoids the singularity: the slices approach the limiting surface $r = 3M/2$, and on the limiting slice $\|K\|_{\op} \le 4\sqrt3/(9M)$ (Table~\ref{tab:noneuclidean_master}), so $K_{ij}K^{ij}$ stays bounded. The known drawback of maximal slicing in numerical evolutions is slice stretching, the growth of $\gamma_{rr}$ near the horizon, not a blow-up of $K_{ij}$.

Motivated by the discrete $\Gamma$-convergence conjecture of Chapter~7 \cite{silvafilho2026minimax}, consider the operator norm $\|K\|_{\op} = \sup_{v} |K(v,v)| / \gamma(v,v)$:
\begin{enumerate}
    \item A minimax slicing would minimize the peak extrinsic curvature,
    \begin{equation}
    \Sigma^* \in \arg\min_{\Sigma \in \mathcal{C}} \operatorname{ess\,sup}_{p \in \Sigma} \|K_p\|_{\op},
    \end{equation}
    over a class $\mathcal{C}$ of slices with prescribed behaviour, for instance slices that agree with a given slice near spatial infinity and cross the future event horizon. Without such a condition the problem is degenerate (Remark~\ref{rem:minimax_slicing}): the infimum is $0$, attained by the slice $t = \text{const}$, which does not enter the black hole. Existence of a minimizer in a class $\mathcal{C}$ of this kind is not addressed here.
    \item \textbf{Strong Hyperbolicity Preservation:} One could try to realize such slicings dynamically by coupling the evolution to the \textbf{Bona-Mass\'o slicing family}:
    \begin{equation}
    \partial_t \alpha - \beta^k \partial_k \alpha = -\alpha^2 f(\alpha) \operatorname{tr}(K)
    \end{equation}
    where $f(\alpha) > 0$. For this family the gauge characteristic speeds, $\alpha\sqrt{f(\alpha)\,\gamma^{xx}}$ along the $x$ direction, are real, a standard ingredient of strong hyperbolicity of BSSN-type systems \cite{gourgoulhon2012}. \emph{We conjecture} that adding the barrier constraint $\|K\|_{\op} \le \kappa^*$ as a limiter on $\partial_t\alpha$ preserves this property and freezes the lapse near the singularity. This has not been shown.
    \item \emph{Conjecturally}, a barrier parameter $\mu \ge r(p)^{-2}$, with $r$ the areal radius, in a discretization of the type conjectured to $\Gamma$-converge in Chapter~7 yields a uniform clearance from the singularity $r = 0$,
    \begin{equation}
    \min_{p \in \Sigma^*} r(p) \ge r_0 > 0 .
    \end{equation}
    The $\Gamma$-convergence statement of Chapter~7 is itself a conjecture, and it concerns Euclidean obstacle problems, so even if true it would not by itself imply this.
\end{enumerate}

\subsection{Curvature Thermodynamics and the Gibbons-Hawking-York Boundary Action}
\label{subsec:curvature_thermodynamics}
A curvature plateau $\kappa(s) \equiv K^*$ on a contact locus $\mathcal{S}$ bears a formal resemblance to boundary terms of the gravitational path integral, which we record here without claiming more.
In quantum gravity, the Euclidean Einstein-Hilbert action with Gibbons-Hawking-York (GHY) boundary term is:
\begin{equation}
I_{\mathrm{grav}}[g] = -\frac{1}{16\pi G} \int_{\mathcal{M}} (R - 2\Lambda) \sqrt{g} \, d^4x - \frac{1}{8\pi G} \oint_{\partial\mathcal{M}} K \sqrt{h} \, d^3x
\end{equation}

\begin{remark}[GHY Boundary Action Minimization and Curvature Equipartition]
\label{thm:curvature_equipartition}
Let $\partial\mathcal{M}$ be an obstacle contact interface. We observe a structural resemblance between minimizing the operator norm $\|K\|_{\op} \le K^*$ and bounding the GHY boundary term:
\begin{equation}
|I_{\mathrm{GHY}}| \le \frac{3}{8\pi G} K^* \operatorname{Vol}_3(\partial\mathcal{M}),
\end{equation}
which follows from $|K| = |\operatorname{tr} K| \le 3\|K\|_{\op} \le 3K^*$ on the three-dimensional boundary. We do not prove any statement about path-integral minimization; we only note that the classical boundary action is controlled by the minimax curvature.
\end{remark}

\begin{remark}[Semi-Classical Limitations and Quantum Gravity]
This boundary-action formulation is \textbf{semi-classical}. Our $L^\infty$ minimax functional operates on smooth Lorentzian manifolds. The Euclidean computation of Gibbons and Hawking, in which the GHY term yields the black-hole entropy, is reliable for black holes much larger than $\ell_P$; at the Planck scale the metric itself fluctuates, the semiclassical saddle-point evaluation is no longer controlled, and the classical curvature bounds used here lose their meaning. They should be read as statements about classical or mean-field geometries.
\end{remark}

\subsection{Gravitational Wave Deflection Rings and Strong Lensing (LIGO/LISA)}
\label{subsec:lensing_ligo_lisa}
In the geometric-optics limit, high-frequency gravitational wave packets propagate along null geodesics. The minimax framework of this chapter does not apply to them as it stands.

\begin{remark}[Null curves]
\label{thm:gw_lensing}
For a null curve with tangent $k$, the tangent line lies in its own orthogonal complement, $k \in k^\perp$, so there is no splitting $T\mathcal{M} = T\gamma \oplus N\gamma$ and no second fundamental form in the sense of Definition~\ref{def:opnorm}; the arc length vanishes, and an affine parameter is defined only up to affine rescaling. Moreover, an affinely parametrized null geodesic has $\bar\nabla_k k = 0$, so any curvature built from $\bar\nabla_k k$ vanishes identically for the paths along which gravitational waves travel. A minimax analysis of lensing, for instance of the relativistic images near the photon sphere $r = 3M$, would have to use a different curvature, such as the geodesic curvature of the projection of the ray into a static slice with the optical metric. We do not pursue this here.
\end{remark}

\subsection{The Big Bounce: Loop Quantum Cosmology, Conformal Factors, and a Planck Density Bound}

When extending the minimax curvature bound to cosmological singularities (such as the Big Bang), the extrinsic shear limit operates as an effective geometric cutoff that mirrors the phenomenology of Loop Quantum Cosmology (LQC) \cite{ashtekar2006quantum}.

\begin{remark}[Minimax Saturation and LQC Holonomy Equivalence]
In a Friedmann-Lema\^{i}tre-Robertson-Walker (FLRW) background, the extrinsic curvature trace $K = -3H$ (where $H = \dot{a}/a$). Enforcing the universal Planck-scale minimax saturation $\|\II_g\|_{\mathrm{op}} \le \ell_P^{-1}$ bounds the Hubble rate $H^2 \le \mathcal{O}(\ell_P^{-2})$. We do not derive this from full Loop Quantum Cosmology holonomy corrections, but observe that this geometric constraint structurally parallels the modified effective Friedmann equation of LQC:
\begin{equation}
H^2 = \frac{8\pi G}{3} \rho \left(1 - \frac{\rho}{\rho_{\mathrm{crit}}}\right)
\end{equation}
where the density peaks at a critical value $\rho_{\mathrm{crit}} \sim \rho_{\mathrm{Pl}}$, replacing the initial singularity with a smooth Big Bounce.
\end{remark}

In SI units the FLRW slices have $K_{ij} = -(H/c)\gamma_{ij}$, so the cutoff $\|K\|_{\op} \le \ell_P^{-1}$ reads $|H| \le c/\ell_P$. For spatially flat or open FLRW ($k \le 0$) with $\Lambda \ge 0$, the Friedmann equation $H^2 = 8\pi G\rho/3 - kc^2/a^2 + \Lambda c^2/3$ gives $H^2 \ge 8\pi G\rho/3$, with equality for $k = 0$, $\Lambda = 0$, so the cutoff bounds the energy density,
\begin{equation}
\rho c^2 \le \frac{3c^4}{8\pi G_N \ell_P^2} = \frac{3}{8\pi}\,\frac{c^7}{\hbar G_N^2} \approx 5.5 \times 10^{112}\ \mathrm{J\,m^{-3}},
\end{equation}
a fraction $3/(8\pi)$ of the Planck energy density $c^7/(\hbar G_N^2) \approx 4.63 \times 10^{113}$\,J\,m$^{-3}$ (numerically equal to the Planck pressure in Pa). For closed FLRW ($k = +1$) the term $-kc^2/a^2$ can offset $8\pi G\rho/3$, and $|H|$ alone bounds nothing. No bound on the pressure follows: in FLRW $p = -(2\dot H + 3H^2)c^2/(8\pi G_N)$ involves $\dot H$, which an $L^\infty$ bound on $K_{ij}$ at each instant does not control.

\begin{remark}[Conformal Factor and Lorentzian Continuation]
For a conformally rescaled metric $\tilde g = \Omega^2 g$, the principal curvatures of a hypersurface transform as $\tilde\kappa_i = \kappa_i/\Omega + \tilde{\mathbf{n}}(\ln\Omega)$, up to the orientation sign, where $\tilde{\mathbf{n}} = \Omega^{-1}\mathbf{n}$ is the $\tilde g$-unit normal. A cutoff $\|\tilde\II\|_{\op} \le \ell_P^{-1}$ on the boundary therefore gives $|\tilde{\mathbf{n}}(\ln\Omega)| \le \ell_P^{-1} + \|\II\|_{\op}/\Omega$, or, with the $g$-unit normal $\mathbf{n}$, $|\mathbf{n}(\ln\Omega)| \le \Omega\,\ell_P^{-1} + \|\II\|_{\op}$. This limits, but does not cure, the Gibbons--Hawking--Perry conformal-factor problem of Euclidean quantum gravity, since the global mode $\Omega \to \infty$ is unaffected. At a bounce the background extrinsic curvature vanishes, $K_{ij}^{(0)} = 0$, so the slice is a moment of time symmetry, and perturbations $(\delta h_{ij}, \delta K_{ij})$ can be evolved in real Lorentzian time by the ADM system. Whether this provides a consistent replacement for the analytic continuation used in the Hartle--Hawking construction is not addressed here.
\end{remark}

\section{Universal Master Classification Table in Non-Euclidean Geometries}

\begin{table}[htbp]
\centering
\scriptsize
\renewcommand{\arraystretch}{1.3}
\begin{tabularx}{\textwidth}{@{}l l c c >{\raggedright\arraybackslash}X >{\raggedright\arraybackslash}X >{\raggedright\arraybackslash}p{2.3cm}@{}}
\toprule
Ambient Space $(N^n, g)$ & Object $M^k$ & $k$ & codim. & Candidate $M^k$ & Physical / Relativistic Context & $\|\II\|_{\op}$ of candidate \\
\midrule
$\Sph^n(+c^2)$ & Curve & $1$ & $n-1$ & Small Circle & Geodesic turning on Sphere & $c \cot(c R)$ \\
$\Sph^2(+c^2)$ & Curve & $1$ & $1$ & Geodesic semicircle, radius $w/2$ & Optimal U-turn in a tube, $w < \pi/c$ (Thm.~\ref{prop:space_form_uturn}) & $c \cot(cw/2)$ \\
$\Sph^3(+1)$ & Surface & $2$ & $1$ & Clifford Torus $S^1 \times S^1$ & Minimal flat surface in $\Sph^3$ & $1$ \\
$\Hyp^2(-c^2)$ & Curve & $1$ & $1$ & Geodesic semicircle, radius $w/2$ & Optimal U-turn in a tube (Thm.~\ref{prop:space_form_uturn}) & $c \coth(cw/2)$ \\
$\Hyp^n(-c^2)$ & Hypersurface & $n-1$ & $1$ & Horosphere & Flat intrinsic boundary & $c$ \\
$\Hyp^n(-c^2)$, $n \ge 3$ & Hypersurface & $n-1$ & $1$ & Tube of radius $d$ about a geodesic & Equidistant from a line & $c\coth(cd)$ \\
$\R^{1,3}$ (Minkowski) & Spacelike & $3$ & $1$ & Hyperboloid $t^2 - |\mathbf{x}|^2 = \tau_0^2$ & Milne slice (constant proper time from the origin) & $\frac{1}{\tau_0}$ \\
Schwarzschild & Spacelike & $3$ & $1$ & Limiting maximal slice ($K = 0$) \cite{estabrook1973} & Singularity-avoiding ADM & $|K^r{}_r| = \frac{3\sqrt{3} M^2}{2 r^3}$ \\
Reissner--Nordstr\"om & 2-Surface & $2$ & $2$ & Horizon in PG slice (Prop.~\ref{prop:rn_horizon}) & MOTS, $\theta_l = 0$ & $\frac{1}{r_+}$ (slicing-dependent) \\
Morris-Thorne & 2-Surface & $2$ & $2$ & Throat in static slice (Prop.~\ref{prop:wormhole_throat}) & Minimal surface; NEC violated & $0$ (intrinsic $K_{\mathrm{Gauss}} = r_0^{-2}$) \\
$\text{FLRW Cosmology}$ & Hypersurface & $3$ & $1$ & Homogeneous spatial slice & Cosmological expansion & $|H(t)|$ (trace $3H$) \\
\bottomrule
\end{tabularx}
\caption{Explicit candidates in non-Euclidean and spacetime geometries, with their exact peak extrinsic curvature. Except for the U-turn rows (Theorem~\ref{prop:space_form_uturn}), these are examples and upper bounds, not proven minimizers; where a proposition is cited, it computes the curvature. For the Schwarzschild limiting slice, $|K^r{}_r|$ attains its maximum $4\sqrt{3}/(9M)$ at $r = 3M/2$.}
\label{tab:noneuclidean_master}
\end{table}

\section{Regularity Invariance in Riemannian Backgrounds}

\begin{conjecture}[Non-Euclidean Regularity Invariance]
\label{thm:noneuclidean_invariance}
Let $(N^n, g)$ be a smooth Riemannian manifold, $\Omega \subset N$ a bounded domain with $C^{1,1}$ boundary, $\Sigma^{k-1} \subset \partial\Omega$ a compact $C^\infty$ submanifold and $V > 0$. Then, for every $r \ge 2$,
\begin{equation}
\inf_{M \in \mathcal{A}_r(N, g, \Omega, \Sigma, V)} \|\II_M\|_{L^\infty} = \inf_{M \in \mathcal{A}_{1,1}(N, g, \Omega, \Sigma, V)} \|\II_M\|_{L^\infty}.
\end{equation}
\end{conjecture}

For $N = \R^n$ this is the regularity-invariance conjecture of Chapter~7. The inequality $\ge$ is immediate from $\mathcal{A}_r \subset \mathcal{A}_{1,1}$. The analogue for timelike curves in a Lorentzian manifold, which Remark~\ref{rem:worldline_regularity} would need, is not covered by this statement.

\begin{remark}[Status]
A naive proof with the cutoff mollifier $h_\epsilon = \chi h + (1-\chi)(h * \rho_\epsilon)$, where $\chi$ switches on a layer of width $\epsilon^2$, fails: the second derivatives of $h_\epsilon$ contain $\chi''(h - h*\rho_\epsilon)$ and $2\chi'\,(Dh - D(h*\rho_\epsilon))$, of orders $\epsilon^{-4}\cdot\epsilon^2$ and $\epsilon^{-2}\cdot\epsilon$, which are not bounded. The Euclidean case is itself open (Chapter~7). We do not know whether the curved case reduces to it. The Euclidean conjecture is a statement about global infima with boundary, volume and obstacle constraints, not a local approximation lemma that could be transported chart by chart, and in a chart the Christoffel symbols enter the second fundamental form, so they change the value of the functional being minimized.
\end{remark}

% Shared declarations block, \input by every chapter before its bibliography.
% Keep in sync with the "Declarations" section of master_book_unified_quantum_gravity.tex.
\section*{Declarations}

\noindent\textbf{Funding.} This research was financed in part by the Coordena\c{c}\~ao de Aperfei\c{c}oamento de Pessoal de N\'ivel Superior -- Brasil (CAPES) -- Finance Code 001.

\medskip\noindent\textbf{Competing interests.} The author declares no competing interests.

\medskip\noindent\textbf{Use of generative AI tools.} Large language model assistants (Anthropic Claude; Google Gemini, used through the Antigravity agent environment) were used extensively in preparing this work: drafting and editing text and \LaTeX{} source; proposing, expanding and revising derivations and proofs under the author's direction; writing the Python verification scripts and the \textsc{Lean 4} files; searching and checking the literature and references; and adversarial auditing of proofs, numerical claims and code. These audits found errors, some of them originally introduced with AI assistance. The errors and their corrections are recorded in the repository file \texttt{WORKPLAN.md}. AI tools are not authors. The author chose the research questions and the overall framework, reviewed the material, and is responsible for the content, including any remaining errors.

\medskip\noindent\textbf{Verification status.} Results labelled Theorem, Proposition or Lemma are proved in the text or cited as classical; statements labelled Conjecture are not proved. The numerical scripts compare formulas of the text with independent computations and are checks, not proofs. The \textsc{Lean 4} modules in \texttt{formal\_proofs\_book/Book} are a naming skeleton and do not verify the mathematics of this chapter.

\medskip\noindent\textbf{Code and data availability.} Sources, numerical scripts and the audit log are available at
\begin{center}\url{https://github.com/reinaldomsilvafilho-netizen/quantum-gravity-lean4}.\end{center}
The monograph is archived on Zenodo, \href{https://doi.org/10.5281/zenodo.22290043}{doi:10.5281/zenodo.22290043}, under the CC~BY~4.0 license.


\begingroup\raggedright
\begin{thebibliography}{99}

\bibitem{silvafilho2026minimax} R.~M. Silva-Filho, \emph{Minimax-Flat $k$-Submanifolds in Obstacle Environments: Variational Theory, Constructive Heuristics, and Explicit Candidates}, Chapter~7 of \emph{Geometry, Tensors, and Quantum Gravity}, Universidade Federal de Lavras, 2026.

\bibitem{adm1962}
R.~Arnowitt, S.~Deser, and C.~W.~Misner,
\emph{The dynamics of general relativity},
in Gravitation: An Introduction to Current Research, L. Witten ed., Wiley, New York (1962), 227--265., \href{https://doi.org/10.1007/s10714-008-0661-1}{doi:10.1007/s10714-008-0661-1}

\bibitem{israel1966}
W.~Israel,
\emph{Singular hypersurfaces and thin shells in general relativity},
Nuovo Cimento B \textbf{44} (1966), 434--449., \href{https://doi.org/10.1007/BF02710419}{doi:10.1007/BF02710419}

\bibitem{morristhorne1988}
M.~S.~Morris and K.~S.~Thorne,
\emph{Wormholes in spacetime and their use for interstellar travel},
Amer. J. Phys. \textbf{56} (1988), 395--412., \href{https://doi.org/10.1119/1.15620}{doi:10.1119/1.15620}

\bibitem{raychaudhuri1955}
A.~Raychaudhuri,
\emph{Relativistic cosmology I},
Phys. Rev. \textbf{98} (1955), 1123--1126., \href{https://doi.org/10.1103/PhysRev.98.1123}{doi:10.1103/PhysRev.98.1123}

\bibitem{hawking1992}
S.~W.~Hawking,
\emph{Chronology protection conjecture},
Phys. Rev. D \textbf{46} (1992), 603--611., \href{https://doi.org/10.1103/PhysRevD.46.603}{doi:10.1103/PhysRevD.46.603}

\bibitem{oneill1983}
B.~O'Neill,
\emph{Semi-Riemannian Geometry with Applications to Relativity},
Academic Press, New York, 1983.

\bibitem{gourgoulhon2012}
E.~Gourgoulhon,
\emph{3+1 Formalism in General Relativity},
Lecture Notes in Physics, vol.~846, Springer, 2012, \href{https://doi.org/10.1007/978-3-642-24525-1}{doi:10.1007/978-3-642-24525-1}

\bibitem{mars1993}
M.~Mars and J.~M.~M. Senovilla,
\emph{Geometry of general hypersurfaces in spacetime: junction conditions},
Class. Quantum Grav. \textbf{10} (1993), 1865--1897, \href{https://doi.org/10.1088/0264-9381/10/9/026}{doi:10.1088/0264-9381/10/9/026}

\bibitem{estabrook1973}
F.~Estabrook, H.~Wahlquist, S.~Christensen, B.~DeWitt, L.~Smarr, and E.~Tsiang,
\emph{Maximally slicing a black hole},
Phys. Rev. D \textbf{7} (1973), 2814--2817, \href{https://doi.org/10.1103/PhysRevD.7.2814}{doi:10.1103/PhysRevD.7.2814}

\bibitem{ashtekar2006quantum}
A.~Ashtekar, T.~Pawlowski, and P.~Singh,
\emph{Quantum Nature of the Big Bang: Improved dynamics},
Phys. Rev. D \textbf{74} (2006), 084003., \href{https://doi.org/10.1103/PhysRevD.74.084003}{doi:10.1103/PhysRevD.74.084003}

\end{thebibliography}
\endgroup

\end{document}
